\section{Fazit und Ausblick auf FFT}
\label{sec:implementationconclusion}

Im Rahmen dieser Seminararbeit wurde erläutert, was die Convolution ist
und wie sie mathematisch auf diskrete Signale angewendet wird.
Dabei wurden anhand der Beispiele FIR-Filter und Convolution Reverb gezeigt, welche die Wirkung Convolution
auf ein Eingangssignal zeigen kann, sowie verschiedene Arten der Umsetzung einer Convolution genannt.
Es wurde auch erläutert, dass es dafür neben dem diskreten Eingangssignal auch eine
diskrete Impulsantwort braucht und welche verschiedenen Techniken es gibt, um diese aus
einem System zu ermitteln.

Wie am Experiment in der vorherigen Sektion \ref{sec:implementation} zu sehen, ist eine
direkte Implementierung der diskreten Convolution \ref{eq:directConv} in Audiosoftware und erst recht
in Echtzeitanwendungen, aufgrund der sehr hohen Laufzeiten, nicht sinnvoll.
Die einzige Möglichkeit der effizienten Anwendung der Convolution liegt deshalb in
der Verwendung der FFT, um die diskreten Signale in ihre Frequenzdomäne zu bringen, sodass
die Convolution in einer einfachen Multiplikation durcheführt werden kann.
Erst dieses Vorgehen macht eine Nutzung der Convolution in der Praxis effizient nutzbar.

